\documentclass[11pt]{amsart}
\usepackage{amssymb}
\usepackage[margin=1.15in]{geometry}
\usepackage[hidelinks]{hyperref}
\newtheorem{theorem}{Theorem}
\newtheorem{corollary}[theorem]{Corollary}
\newcommand{\content}{\mathrm{c}}
\newcommand{\code}[1]{\texttt{#1}}
\title{The direct Dedekind--Mertens proof and its Unit formalisation}
\date{October 25, 2026}
% stamp on every page: file name and date/time of compilation
\newcount\stamphour \newcount\stampmin
\stamphour=\time \divide\stamphour by 60
\stampmin=\time \multiply\stamphour by 60 \advance\stampmin by -\stamphour \divide\stamphour by 60
\edef\stamptime{\the\year-\ifnum\month<10 0\fi\the\month-\ifnum\day<10 0\fi\the\day\space
 \ifnum\stamphour<10 0\fi\the\stamphour:\ifnum\stampmin<10 0\fi\the\stampmin}
\AddToHook{shipout/foreground}{\put(0.5\paperwidth,-\dimexpr\paperheight-0.45in\relax)%
 {\makebox[0pt]{\scriptsize\ttfamily\jobname.pdf\quad\stamptime}}}
\begin{document}
\maketitle

This note presents the argument in the supplied \emph{A direct proof of
the Dedekind-Mertens Lemma}, dated April 5, 2006
(\code{mertens.pdf}), and describes its checked formalisation in Unit.
The proof works over every commutative ring, without a noetherian
hypothesis.

\section{The subgroup theorem}

Let $A$ be a commutative ring and let
\[
 P=\sum_{i=0}^n f_iX^i,\qquad
 Q=\sum_{j=0}^m g_jX^j,\qquad
 PQ=\sum_k h_kX^k.
\]
The bounds $n,m$ are formal degrees: the displayed leading coefficients
need not be nonzero. Let $F,G,H$ be the additive subgroups of $A$
generated by the coefficients of $P,Q,PQ$, respectively.
For additive subgroups $L,M$, write $LM$ for the additive subgroup
generated by all $ab$ with $a\in L$, $b\in M$. Set $F^0=\mathbb Z1$.
For $a\in A$, the notation $aL$ denotes the subgroup of multiples
$a\ell$ with $\ell\in L$.

\begin{theorem}\label{thm:strong}
One has $F^{m+1}G=F^mH$.
\end{theorem}
\begin{proof}
The inclusion $F^mH\subseteq F^{m+1}G$ follows from the coefficient
formula for $PQ$, which gives $H\subseteq FG$.
We prove the other inclusion by induction on $m$.
For $m=0$, the coefficients of $PQ$ are $f_i g_0$, so $FG=H$.

Suppose $m>0$, set $Q_-=Q-g_mX^m$, and let $G_-$ be the subgroup
generated by $g_0,\ldots,g_{m-1}$. Extend the notation by $f_i=0$
outside $0\leq i\leq n$. Every coefficient of $PQ_-$ is
$h_k-f_{k-m}g_m$, hence lies in $H+g_mF$. The induction hypothesis
applied to $P,Q_-$ gives
\[
 F^mG_-\subseteq F^{m-1}H+g_mF^m.
\]
Writing $K=F^mH$ and multiplying by $f_i$, we obtain
\begin{equation}\label{eq:outer}
 f_iF^mG_-\subseteq K+F^m f_i g_m.
\end{equation}
The coefficient $h_{i+m}$ gives
\[
 f_i g_m
 =h_{i+m}-\sum_{j<m} f_{i+m-j}g_j
 \in H+\sum_{\ell>i}f_\ell G_-.
\]
Together with~\eqref{eq:outer}, this implies
\[
 f_iF^mG_-\subseteq K+\sum_{\ell>i}f_\ell F^mG_-.
\]
Descending induction on $i=n,n-1,\ldots,0$ therefore proves
$f_iF^mG_-\subseteq K$ for every $i$, hence
$F^{m+1}G_-\subseteq K$.

To include the remaining coefficient $g_m$, multiply the same
coefficient identity by $F^m$. The terms
$f_\ell F^mG_-$ already belong to $K$, so
$F^m f_i g_m\subseteq K$ for every $i$. Thus
$F^{m+1}g_m\subseteq K$ as well. Since
$G=G_-+\mathbb Zg_m$, the required inclusion follows.
\end{proof}

\begin{corollary}[Dedekind--Mertens]
Writing $\content(P)$ for the ideal generated by the coefficients of
$P$, one has
\[
 \content(P)^{m+1}\content(Q)=\content(P)^m\content(PQ).
\]
\end{corollary}
\begin{proof}
Take the ideals generated by the two equal subgroups in
Theorem~\ref{thm:strong}. Their generators are respectively
$f_{i_0}\cdots f_{i_m}g_j$ and
$f_{i_1}\cdots f_{i_m}h_k$.
\end{proof}

\begin{corollary}\label{cor:zero}
If $PQ=0$, then
\[
 f_{i_0}\cdots f_{i_m}g_j=0
\]
for every tuple of indices. In particular,
$\content(P)^{m+1}\content(Q)=0$.
\end{corollary}

\section{The coefficient order used by Unit}

The polynomial library stores coefficients with the constant term
first. The formal proof therefore removes the constant term of $Q$
and uses ascending induction. This reverses the coefficient order
in the preceding argument without changing its algebraic content.

Write $Q=g+XQ'$, let $G'$ be the subgroup generated by the
coefficients of $Q'$, and put $H'=\langle\text{coefficients of }PQ'\rangle$.
For $m>0$ the outer induction hypothesis is
$F^mG'\subseteq F^{m-1}H'$. Put $K=F^mH$ and consider
\[
 I(i):\qquad f_iF^mG'\subseteq K.
\]
The coefficient identity is
\[
 h_i=g f_i+(PQ')_{i-1},
\]
where the last term is zero for $i=0$.
If $I(\ell)$ holds for all $\ell<i$, the convolution formula for
$(PQ')_{i-1}$ gives
$F^m(PQ')_{i-1}\subseteq K$. It follows that
\begin{equation}\label{eq:head}
 F^m f_i g\subseteq K.
\end{equation}
To prove $I(i)$, multiply the outer induction inclusion by $f_i$.
Each resulting generator has the form $f_i t(PQ')_j$, where $t$
is a product of $m-1$ coefficients of $P$. Now
\[
 f_i t(PQ')_j=f_i t h_{j+1}-t f_{j+1}f_i g.
\]
The first term belongs to $K$ by definition, and the second belongs
to $K$ by~\eqref{eq:head}, since $tf_{j+1}\in F^m$.
This proves $I(i)$. Primitive recursion supplies the finite list of
earlier statements needed at each index. Finally, $I(i)$ handles
the coefficients of $Q'$, while~\eqref{eq:head} handles $g$.

\section{What has been formalised}

The development is in \code{unit/sag/mertens/}. It checks the entire
univariate subgroup theorem, both inclusions of the ideal formula,
and the coefficient-annihilation assertion of Corollary~\ref{cor:zero}.
The coefficient ring is arbitrary.

Membership is represented by explicit certificates. For a family
$(a_i)$, a subgroup certificate is a finite expression built from
$0$, the $a_i$, addition and negation, together with an equality
identifying its value. Ideal certificates additionally allow arbitrary
ring multiples of the generators. Thus the strong theorem really
uses only integer signed sums. Equality of the generated subgroups
or ideals is expressed by maps on certificates in both directions;
no equality of the certificate types is asserted.

\begin{itemize}
\item \code{AdditiveSpan.seg} defines signed-sum certificates,
their elimination into additive subgroups, and bounded induction.
\item \code{CoefficientProducts.seg} proves the product-coefficient
formula with the bounds required by the inner induction.
\item \code{DedekindMertens.seg} contains the two inductions and
the explicit certificates for $F^{m+1}G\subseteq F^mH$.
\item \code{IdealSpan.seg} defines ideal certificates.
\code{ContentFormula.seg} proves both subgroup inclusions and both
ideal inclusions, using coefficient-product generating families
directly.
\item \code{ZeroProduct.seg} proves coefficient annihilation when
$PQ$ is zero coefficientwise.
\end{itemize}

For a coefficient list $p$ representing $P$, and $Q=\code{cons}\ g\ q$
with $m=\code{length}\ q$, the main interface is
\begin{verbatim}
module D = DedekindMertens[l].Over[A]
D.theorem p q g word j
\end{verbatim}
Here \code{word} has $m+1$ coefficient indices; the result is a
signed-sum certificate for their coefficient product times $g_j$.
The full formulas are available through
\begin{verbatim}
module E = ContentFormula[l].Over[A,p,g,q]
E.subgroup-forward a certificate
E.subgroup-backward a certificate
E.ideal-forward a certificate
E.ideal-backward a certificate
\end{verbatim}

The proof uses Unit's data eliminators and equality rules. It adds
no axiom and requires neither univalence or replacement parameters,
choice, excluded middle, nor a noetherian hypothesis. It uses raw
coefficient lists, so it does not require polynomial quotients or a
decision procedure for degrees. Zero polynomials and zero leading
coefficients are allowed.

From \code{unit/}, the checking command is
\begin{verbatim}
./segc sag/mertens all
\end{verbatim}
The check on October 25, 2026 reported \emph{14 modules, 0 failed},
including imported library modules.

The multivariate content lemma and its localization application in
\code{CAR1.pdf} have not yet been connected to this development.
The separate Jacobian formalisation still assumes its explicit
annihilator-stabilization input; the noetherian reduction remains
outside that formalisation. Those limitations concern the application,
not the Dedekind--Mertens theorem proved here.
\end{document}
